\documentclass{article}
\usepackage{graphicx}
\usepackage{amsmath}
\usepackage{amssymb}
\usepackage{amsfonts}
\usepackage{hyperref}
\usepackage{exercises}
\usepackage{indentfirst}
\usepackage{circledsteps}
\usepackage{amsthm}
\usepackage{csquotes}
\usepackage{yhmath}
\usepackage{comment}
\usepackage{placeins}
\usepackage{geometry}
\usepackage{hyperref}
\usepackage{tikz}
\newcommand*{\oldneg}{\mathord{\sim}}
\newcommand{\Z}{\mathbb{Z}}
\newcommand{\N}{\mathbb{N}}
\let\oldemptyset\emptyset

\theoremstyle{definition}
\newtheorem{definition}{Definition}
\newtheorem{theorem}{Theorem}[section]
\newtheorem{corollary}{Corollary}[theorem]
\newtheorem{lemma}[theorem]{Lemma}

\setlength{\parindent}{.5in}
\geometry{margin=0.5in}
\setlength{\parskip}{12pt}
\title{Chapter 7 Notes}
\begin{document}
\date{}
\maketitle

\section*{Section 1}
In the Remarks are page 91, it is said that there is an \textquote{easy induction} showing that $R_{n+1}$ has $2^n$ elements. However, this is clearly 
not the case. For, as remakred, $R_3$ has 4 elements. Then $R_4$ has $2^4=16$ elements, which is not equal to $2^3=8$. 

The proof of Proposition 1.3 is problematic. It uses Cantor's Theorem, though this hasn't even been stated (let alone proven) in the context of 
Zermelo-Frankel Universes. Proposition 1.3 can also be proven without using Cantor's Theorem. In particular, we first show that $R_\alpha\not\in R_\alpha$.

\begin{proof}
By transfinite induction. Clearly this holds for 0. Suppose it holds for some $\alpha$. Then, if $R_{\alpha+1}\in R_{\alpha+1}$, by definition 
of $R_{\alpha+1}$, $R_{\alpha+1}$ would be a subset of $R_\alpha$. But $R_\alpha\in R_{\alpha+1}$. Thus this would mean that 
$R_\alpha\in R_\alpha$, contradicting the induction assumption. Hence $R_{\alpha+1}\notin R_{\alpha+1}$. On the other hand, suppose it held for all 
ordinals less than some limit ordinal $\lambda$. Then, by definition, $R_\lambda=\cup_{\alpha<\lambda}R_\alpha$. Suppose 
$R_\lambda\in R_\lambda$. Then, for some $\alpha<\lambda$, $R_\lambda\in R_\alpha$. But $R_{\alpha+1}\subset R_\lambda$, and 
$R_\alpha\in R_{\alpha+1}$, meaning $R_\alpha\in R_\lambda$. Thus, because $R_\alpha$ is transitive, $R_\alpha\in R_\alpha$, contradicting the 
induction assumption. Thus $R_\lambda\notin R_\lambda$, completing the transfinite induction.
\end{proof}

Then, in the proof of Proposition 1.3, $R_\alpha\ne R_{\alpha+1}$ follows for the following reason: We know that 
$R_\alpha\in R_{\alpha+1}$ and $R_\alpha\notin R_\alpha$ (by the above result), meaning 
$R_\alpha\ne R_{\alpha+1}$.

\section*{Section 2}
Note that, at a minimum, the $R_\alpha$ sequence discussed in this section is an $On$-sequence of subsets of $V$ (and of course, all $On$-sequences for any 
particular class are also $On$-sequences for $V$ itself).

Note that the proof of $O_6$ uses Theorem 1.18 of Chapter 5.

\end{document}
